\documentclass[stage=proposal,campus=shenzhen,degree-level=master,lang=en,fontset=fandol]{hithesis}
\hitsetup{top-level-heading-pagebreak-before=true}
\begin{document}
\section{Alpha heading}
BEFORE FLOAT. Some text before the deferred figure.
\begin{figure}[p]
\centering\rule{80mm}{35mm}
\caption{DEFERRED FIGURE}
\end{figure}
\section{Beta heading}
AFTER FLOAT. This heading must follow the deferred figure.
\sectionbreak[continuous]
CONTINUOUS MARKER. The section counter of the report is still \thesection.
\sectionbreak[next]
NEXT MARKER. Page \thepage.
\sectionbreak[odd]
ODD ONE. Page \thepage.
\sectionbreak[odd]
ODD TWO. Page \thepage.
\sectionbreak[even]
EVEN ONE. Page \thepage.
\sectionbreak[even]
EVEN TWO. Page \thepage.
\sectionbreak
DEFAULT NEXT. Page \thepage.
\sectionbreak[next,page-number-format={number-format=lower-roman,start-at=1}]
ROMAN PAGE: \thepage.
\sectionbreak[continuous,columns={number-of-columns=2,spacing=8mm,line-between=true}]
\noindent COLUMN A. The first column contains a short paragraph. This text checks
that a continuous section can enter a two-column layout without forcing a page.
\par\noindent COLUMN B. The second paragraph remains in the same section and is
balanced when the following continuous section returns to a single column.
\sectionbreak[continuous,columns={number-of-columns=1}]
SINGLE AGAIN. The report section number is still \thesection.
\sectionbreak[next,page-setup={paper={width=210mm,height=297mm},
  margins={left=4cm,right=4cm,top=4cm,bottom=3cm},
  layout={from-edge={header=2cm,footer=2cm}}}]
GEOMETRY MARKER.\par
\noindent\rule{\textwidth}{1bp}
\end{document}
